\section{Choosing calculus tools}

The final skill of this chapter is choosing the right tool. In real problems, no one announces in advance whether the solution requires a derivative, an integral, optimisation, or a differential equation. The wording of the problem must be read carefully.

\subsection*{Questions that suggest derivatives}

Use derivatives when the problem asks about local change, slope, rate, growth, decrease, velocity, acceleration, or sensitivity.

Typical phrases include:
\begin{itemize}
\item instantaneous velocity,
\item rate of change,
\item increasing or decreasing,
\item maximum or minimum,
\item fastest increase or decrease,
\item tangent slope.
\end{itemize}

For example, if a concentration is given by \(C(t)\), then the time of maximum concentration is found by studying \(C'(t)\). If the question asks when the concentration decreases most rapidly, then the derivative and second derivative may both be relevant.

\subsection*{Questions that suggest integrals}

Use integrals when the problem asks for accumulation, total amount, area, volume, work, distance from velocity, or average value.

Typical phrases include:
\begin{itemize}
\item total distance,
\item accumulated amount,
\item area under a graph,
\item work done by a variable force,
\item volume of a solid of revolution,
\item average value over an interval.
\end{itemize}

For example, if a force depends on position, work is computed by integrating force over displacement. If a flow rate depends on time, total flow is computed by integrating the rate over time.

\subsection*{Questions that suggest optimisation}

Optimisation problems ask for a best possible value under constraints. They often use derivatives, but the model-building step comes first.

Typical phrases include:
\begin{itemize}
\item maximise area,
\item minimise cost,
\item use the least material,
\item find the best dimensions,
\item largest volume,
\item shortest time or distance.
\end{itemize}

The important step is to build an objective function and identify the allowed domain.

\subsection*{Questions that suggest differential models}

A differential model appears when the problem describes a law of change rather than the quantity itself.

Typical phrases include:
\begin{itemize}
\item the rate is proportional to the current amount,
\item acceleration is constant,
\item velocity is the derivative of position,
\item the quantity changes according to its present value.
\end{itemize}

Such problems often require integration or solving a simple differential equation with an initial condition.

\subsection*{A final comparison}

\begin{center}
\begin{tabular}{lll}
\toprule
Problem type & Mathematical object & Common tool \\
\midrule
Local change & slope or rate & derivative \\
Optimisation & best value of a function & derivative and domain check \\
Accumulation & total from small contributions & definite integral \\
Motion & position, velocity, acceleration & derivatives and integrals \\
Work & force over distance & definite integral \\
Growth model & rate law & differential equation \\
Geometry & area, volume, length & definite integral \\
\bottomrule
\end{tabular}
\end{center}

This table should be read as a guide. Some problems use more than one tool. For example, an optimisation problem may require geometry to build a function, derivatives to find a critical point, and endpoint checks to confirm the result.

\begin{summarybox}
When facing an application problem, ask:
\begin{itemize}
\item What is being asked: rate, total, best value, or model?
\item What quantity should become a function?
\item What is the independent variable?
\item What restrictions or units are present?
\item Which calculus tool matches the question?
\item How should the mathematical result be interpreted in context?
\end{itemize}
\end{summarybox}

This chapter closes the calculus part of the course by showing that calculus is a language for reasoning. Derivatives and integrals are not only techniques. They are tools for turning change and accumulation into mathematical understanding.